\documentclass[12pt,a4paper]{article}
\usepackage[brazilian]{babel}
\usepackage{amsmath,amssymb,amsthm}
\usepackage{geometry}
\usepackage{enumitem}
\usepackage{hyperref}
\usepackage{xcolor}
\usepackage{graphicx}
\usepackage{fancyhdr}
\usepackage{tcolorbox}
\usepackage{totpages}
\usepackage{titlesec}

% Margens ajustadas para cabeçalho institucional
\geometry{top=3.5cm, bottom=2.5cm, left=2.0cm, right=2.0cm, headheight=2.6cm}

% Paleta de Cores (Mesma da Apostila)
\definecolor{primary}{RGB}{10, 45, 90}       % Azul Escuro Institucional
\definecolor{secondary}{RGB}{25, 110, 180}   % Azul Destaque
\definecolor{bglight}{RGB}{245, 247, 250}    % Cinza Claro para Caixas
\definecolor{borderlight}{RGB}{210, 220, 230}% Borda Suave
\definecolor{correctgreen}{RGB}{40, 140, 60}  % Verde para Gabarito

% Configuração das Caixas (tcolorbox)
\tcbuselibrary{skins, breakable}

% Contador de Questões
\newcounter{numquestao}
\setcounter{numquestao}{0}

% Ambiente de Questão Estilizado
\newenvironment{questao}[2][]{%
  \refstepcounter{numquestao}%
  \vspace{0.3cm}%
  \begin{tcolorbox}[
    enhanced,
    colback=white,
    colframe=secondary,
    arc=2pt,
    boxrule=0.8pt,
    leftrule=4pt, % Linha lateral mais espessa para destaque
    title={\textbf{Questão \thenumquestao} \ifstrempty{#2}{}{--- #2} \hfill \small\textnormal{#1}},
    fonttitle=\sffamily\bfseries\color{white},
    coltitle=white,
    colbacktitle=secondary,
    breakable
  ]%
}{%
  \end{tcolorbox}%
}

% Ambiente de Gabarito / Resolução
\newenvironment{gabarito}{%
  \vspace{0.1cm}%
  \begin{tcolorbox}[
    enhanced,
    colback=bglight,
    colframe=correctgreen,
    arc=2pt,
    boxrule=0.6pt,
    title={\small\sffamily\bfseries Resolução / Gabarito},
    fonttitle=\sffamily\color{white},
    colbacktitle=correctgreen,
    breakable
  ]%
  \small
}{%
  \end{tcolorbox}%
}

% Comando para gerar linhas de resposta/rascunho
\newcommand{\espacoresposta}[1]{%
  \vspace{#1}%
}

% Configuração de Cabeçalho e Rodapé
\pagestyle{fancy}
\fancyhf{}

% Cabeçalho com Logo e Disciplina
\fancyhead[L]{%
  \includegraphics[height=2cm]{Imagens/logocaruaru.png}%
}
\fancyhead[R]{%
  \small\sffamily
  \textbf{Instituto Federal de Pernambuco -- Campus Caruaru} \\
  \textcolor{secondary}{Disciplina:} Cálculo I \\
  \textcolor{secondary}{Prof.:} Cleibson José da Silva
}

% Rodapé
\fancyfoot[L]{\small\sffamily\color{gray}Lista 7 -- Volumes}
\fancyfoot[R]{\small\sffamily\color{gray}Página \thepage\ de \ref{TotPages}}

% Linhas divisórias
\renewcommand{\headrulewidth}{0.8pt}
\renewcommand{\headrule}{{\color{primary}\hrule width\headwidth height \headrulewidth}}
\renewcommand{\footrulewidth}{0.4pt}

\hypersetup{
    colorlinks=true,
    linkcolor=secondary,
    urlcolor=secondary,
    citecolor=primary
}

\begin{document}

% Cabeçalho de Identificação do Estudante
\begin{tcolorbox}[colback=bglight, colframe=primary, boxrule=1pt, arc=3pt]
  \sffamily\small
  \begin{tabular}{@{}p{0.65\textwidth} p{0.3\textwidth}@{}}
    \textbf{Estudante:} \hrulefill & \textbf{Data:} \hrulefill \\[0.25cm]
    \textbf{Curso/Turma:} \hrulefill & \textbf{Nota/Visto:} \hrulefill
  \end{tabular}
\end{tcolorbox}

% Título Principal da Lista
\begin{center}
  {\LARGE\bfseries\color{primary} Lista 7: Volumes} \\[0.1cm]
  {\small\sffamily\color{gray} Volumes de Sólidos de Revolução}
\end{center}

\vspace{0.2cm}



\textbf{1--6} Encontre o volume do sólido obtido pela rotação da região limitada pelas curvas dadas em torno das retas especificadas. Esboce a região, o sólido e um disco ou arruela típicos.

\begin{enumerate}
    \item $y = 2 - \frac{1}{2}x, \; y = 0, \; x = 1, \; x = 2; \quad \text{em torno do eixo } x$
    \item $y = e^x, \; y = 0, \; x = 0, \; x = 1; \quad \text{em torno do eixo } x$
    \item $y = 1/x, \; x = 1, \; x = 2, \; y = 0; \quad \text{em torno do eixo } x$
    \item $y = \sqrt{25 - x^2}, \; y = 0, \; x = 2, \; x = 4; \quad \text{em torno do eixo } x$
    \item $x = 2\sqrt{y}, \; x = 0, \; y = 9; \quad \text{em torno do eixo } y$
    \item $y = \ln x, \; y = 1, \; y = 2, \; x = 0; \quad \text{em torno do eixo } y$
\end{enumerate}

\end{document}